\section{The Levi-Civita Connection}

Unless stated otherwise, the constructions below apply to both Riemannian and
pseudo-Riemannian metrics.

\section{The Tangential Connection Revisited}

\begin{proposition}
\label{prop:tangential-covariant-derivative-projection}
\dcref{ch5:5.1}
\uses{prop:adapted-orthonormal-frames,prop:connection-on-tensor-bundles}

Let $M \subseteq \mathbb{R}^n$ be an embedded submanifold, let
$\gamma:I\to M$ be smooth, and let $V$ be a smooth vector field along
$\gamma$ with values in $TM$. Then
\[
D_t^\top V(t)=\pi^\top\bigl(\overline D_tV(t)\bigr)
\]
for every $t\in I$.

\begin{proof}
\sketch Fix $t_0\in I$ and choose an ambient orthonormal frame
$(E_1,\ldots,E_n)$ near $\gamma(t_0)$ such that
$(E_1,\ldots,E_k)$ restricts to a frame for $TM$, where $k=\dim M$.
Writing $V=V^iE_i$ for $1\leq i\leq k$ and using the covariant derivative
formula along a curve gives
\[
\pi^\top(\overline D_tV)
=\sum_{i=1}^k\left(\dot V^iE_i+V^i\pi^\top(\overline\nabla_{\dot\gamma}E_i)\right)
=\sum_{i=1}^k\left(\dot V^iE_i+V^i\nabla^\top_{\dot\gamma}E_i\right)
=D_t^\top V.
\]
\end{proof}
\end{proposition}

\begin{corollary}
\label{cor:geodesic-submanifold-acceleration}
\dcref{ch5:5.2}
\uses{ex:euclidean-connection,prop:tangential-covariant-derivative-projection}

If $M\subseteq\mathbb R^n$ is an embedded submanifold, a smooth curve
$\gamma:I\to M$ is a geodesic for the tangential connection if and only if
$\gamma''(t)\perp T_{\gamma(t)}M$ for every $t\in I$.

\begin{proof}
\sketch The Euclidean connection has zero coefficients in the standard frame,
so $\overline D_t\gamma'=\gamma''$. Apply
Proposition~\ref{prop:tangential-covariant-derivative-projection}.
\end{proof}
\end{corollary}

For pseudo-Euclidean space $(\mathbb R^{r,s},\overline g^{(r,s)})$ and an
embedded nondegenerate submanifold $M$, write
$T_p\mathbb R^{r,s}=T_pM\oplus N_pM$, where
$N_pM=(T_pM)^\perp$, and let $\pi^\top$ denote the corresponding orthogonal
projection. The tangential connection is
\[
\nabla_X^\top Y=\pi^\top(\overline\nabla_{\widetilde X}\widetilde Y),
\]
for ambient extensions $\widetilde X,\widetilde Y$.

\begin{proposition}
\label{prop:geodesic-pseudo-euclidean-submanifold}
\dcref{ch5:5.3}

Suppose $M$ is an embedded Riemannian or pseudo-Riemannian submanifold of
$\mathbb R^{r,s}$. A smooth curve $\gamma:I\to M$ is a geodesic for
$\nabla^\top$ if and only if $\gamma''(t)$ is
$\overline g^{(r,s)}$-orthogonal to $T_{\gamma(t)}M$ for every $t\in I$.
\end{proposition}

\begin{exercise}
\label{exer:prove-geodesic-pseudo-euclidean-submanifold}
\dcref{ch5:5.4}
\uses{prop:geodesic-pseudo-euclidean-submanifold}

Prove Proposition~\ref{prop:geodesic-pseudo-euclidean-submanifold}.
\end{exercise}

\section{Connections on Abstract Riemannian Manifolds}

\section{Metric Connections}

Let $g$ be a Riemannian or pseudo-Riemannian metric on a smooth manifold $M$,
possibly with boundary. A connection $\nabla$ on $TM$ is \emph{compatible
with $g$}, or a \emph{metric connection}, if
\begin{equation}
\nabla_X\langle Y,Z\rangle
=\langle\nabla_XY,Z\rangle+\langle Y,\nabla_XZ\rangle
\tag{5.1}
\end{equation}
for all $X,Y,Z\in\mathfrak X(M)$.

\begin{proposition}[Characterizations of Metric Connections]
\label{prop:characterizations-of-metric-connections}
\dcref{ch5:5.5}

Let $(M,g)$ be a Riemannian or pseudo-Riemannian manifold, possibly with
boundary, and let $\nabla$ be a connection on $TM$. The following are
equivalent:
\begin{enumerate}
\item[(a)] $\nabla$ is compatible with $g$:
\[
\nabla_X\langle Y,Z\rangle
=\langle\nabla_XY,Z\rangle+\langle Y,\nabla_XZ\rangle.
\]
\item[(b)] $g$ is parallel: $\nabla g\equiv0$.
\item[(c)] In every smooth local frame $(E_i)$, the connection and metric
coefficients satisfy
\begin{equation}
\Gamma^l_{ki}g_{lj}+\Gamma^l_{kj}g_{il}=E_k(g_{ij}).
\tag{5.2}
\end{equation}
\item[(d)] For smooth vector fields $V,W$ along any smooth curve $\gamma$,
\begin{equation}
\frac d{dt}\langle V,W\rangle
=\langle D_tV,W\rangle+\langle V,D_tW\rangle.
\tag{5.3}
\end{equation}
\item[(e)] The inner product of two parallel vector fields along a smooth curve
is constant.
\item[(f)] Every parallel transport map along a smooth curve is a linear
isometry.
\item[(g)] Every orthonormal basis at a point of a smooth curve extends to a
parallel orthonormal frame along the curve.
\end{enumerate}

\begin{proof}
\sketch The identity
\[
(\nabla g)(Y,Z,X)
=X(g(Y,Z))-g(\nabla_XY,Z)-g(Y,\nabla_XZ)
\]
proves (a)$\Leftrightarrow$(b), and its frame components
$g_{ij;k}=E_k(g_{ij})-\Gamma^l_{ki}g_{lj}-\Gamma^l_{kj}g_{il}$ prove
(b)$\Leftrightarrow$(c). Expanding $V=V^i\partial_i$ and
$W=W^j\partial_j$ along a curve proves (a)$\Rightarrow$(d); applying (d) to
extendible fields and using the characterization of covariant differentiation
along curves gives the converse.

Condition (d) makes the inner product of parallel fields constant, proving
(e). Parallel extension of two initial vectors then shows that transport
preserves their inner product, proving (f). Transporting an orthonormal basis
gives (g). Finally, in a parallel orthonormal frame, both sides of (5.3) equal
\[
g_{ij}(\dot V^iW^j+V^i\dot W^j),
\]
which proves (g)$\Rightarrow$(d).
\end{proof}
\end{proposition}

\begin{corollary}
\label{cor:metric-connection-norm-geodesic}
\dcref{ch5:5.6}

Let $(M,g)$ be Riemannian or pseudo-Riemannian, possibly with boundary, let
$\nabla$ be a metric connection, and let $\gamma:I\to M$ be smooth.
\begin{enumerate}
\item[(a)] $|\gamma'(t)|$ is constant if and only if
$D_t\gamma'(t)\perp\gamma'(t)$ for every $t\in I$.
\item[(b)] If $\gamma$ is a geodesic, then $|\gamma'(t)|$ is constant.
\end{enumerate}
\end{corollary}

\begin{exercise}
\label{exer:prove-corollary-5-6}
\dcref{ch5:5.7}

Prove Corollary~\ref{cor:metric-connection-norm-geodesic}.
\end{exercise}

\begin{proposition}
\label{prop:tangential-connection-compatible-metric}
\dcref{ch5:5.8}

If $M$ is an embedded Riemannian or pseudo-Riemannian submanifold of
$\mathbb R^n$ or $\mathbb R^{r,s}$, its tangential connection is compatible
with the induced metric.

\begin{proof}
\sketch For ambient extensions $\widetilde X,\widetilde Y,\widetilde Z$ of
$X,Y,Z\in\mathfrak X(M)$, ambient metric compatibility and orthogonal
projection give
\begin{align*}
X\langle Y,Z\rangle
&=\langle\overline\nabla_{\widetilde X}\widetilde Y,\widetilde Z\rangle
 +\langle\widetilde Y,\overline\nabla_{\widetilde X}\widetilde Z\rangle\\
&=\langle\nabla_X^\top Y,Z\rangle+\langle Y,\nabla_X^\top Z\rangle.
\end{align*}
\end{proof}
\end{proposition}

\section{Symmetric Connections}

A connection $\nabla$ on $TM$ is \emph{symmetric} if
\[
\nabla_XY-\nabla_YX=[X,Y]
\qquad(X,Y\in\mathfrak X(M)).
\]
Its torsion tensor is the $(1,2)$-tensor
\[
\tau(X,Y)=\nabla_XY-\nabla_YX-[X,Y].
\]
Thus $\nabla$ is symmetric exactly when $\tau=0$, equivalently when
$\Gamma^k_{ij}=\Gamma^k_{ji}$ in every coordinate frame.

\begin{proposition}
\label{prop:tangential-connection-symmetric}
\dcref{ch5:5.9}
\uses{prop:naturality-of-lie-brackets}

If $M$ is an embedded (pseudo-)Riemannian submanifold of a
(pseudo-)Euclidean space, its tangential connection is symmetric.

\begin{proof}
\sketch Let $\widetilde X,\widetilde Y$ be ambient extensions of
$X,Y\in\mathfrak X(M)$. Naturality of the Lie bracket for the inclusion shows
that $[\widetilde X,\widetilde Y]|_M=[X,Y]$ and is tangent to $M$. Hence
\[
\nabla_X^\top Y-\nabla_Y^\top X
=\pi^\top\bigl([\widetilde X,\widetilde Y]|_M\bigr)=[X,Y].
\]
\end{proof}
\end{proposition}

\begin{theorem}[Fundamental Theorem of Riemannian Geometry]
\label{thm:fundamental-theorem-riemannian-geometry}
\dcref{ch5:5.10}

Let $(M,g)$ be a Riemannian or pseudo-Riemannian manifold, possibly with
boundary. There is a unique symmetric connection $\nabla$ on $TM$ compatible
with $g$. It is the \emph{Levi-Civita connection} of $g$, also called the
\emph{Riemannian connection} when $g$ is positive definite.

\begin{proof}
\sketch Metric compatibility applied to cyclic permutations of $X,Y,Z$, then
combined with symmetry, yields
\begin{align}
\langle\nabla_XY,Z\rangle
&=\tfrac12\bigl(X\langle Y,Z\rangle+Y\langle Z,X\rangle-Z\langle X,Y\rangle
\notag\\
&\qquad-\langle Y,[X,Z]\rangle-\langle Z,[Y,X]\rangle
+\langle X,[Z,Y]\rangle\bigr).
\tag{5.5}
\end{align}
Nondegeneracy of $g$ makes this formula determine $\nabla_XY$, proving
uniqueness.

In a coordinate chart, (5.5) gives
\begin{equation}
\langle\nabla_{\partial_i}\partial_j,\partial_l\rangle
=\frac12(\partial_i\langle\partial_j,\partial_l\rangle
+\partial_j\langle\partial_l,\partial_i\rangle
-\partial_l\langle\partial_i,\partial_j\rangle),
\tag{5.6}
\end{equation}
hence
\begin{equation}
\Gamma^m_{ij}g_{ml}
=\frac12(\partial_ig_{jl}+\partial_jg_{il}-\partial_lg_{ij})
\tag{5.7}
\end{equation}
and
\begin{equation}
\Gamma^k_{ij}
=\frac12g^{kl}(\partial_ig_{jl}+\partial_jg_{il}-\partial_lg_{ij}).
\tag{5.8}
\end{equation}
These coefficients define local connections that agree on overlaps by
uniqueness. They are symmetric in $i,j$, and (5.7) implies
\[
\Gamma^l_{ki}g_{lj}+\Gamma^l_{kj}g_{il}=\partial_kg_{ij},
\]
so Proposition~\ref{prop:characterizations-of-metric-connections} gives metric
compatibility.
\end{proof}
\end{theorem}

\begin{corollary}[Formulas for the Levi-Civita Connection]
\label{cor:formulas-levi-civita-connection}
\dcref{ch5:5.11}
\uses{prob:torsion-tensor,prop:characterizations-of-metric-connections}

Let $(M,g)$ be a Riemannian or pseudo-Riemannian manifold, possibly with
boundary, and let $\nabla$ be its Levi-Civita connection.
\begin{enumerate}
\item[(a)] For smooth vector fields $X,Y,Z$,
\begin{equation}
\langle\nabla_XY,Z\rangle
=\frac12\bigl(X\langle Y,Z\rangle+Y\langle Z,X\rangle-Z\langle X,Y\rangle
-\langle Y,[X,Z]\rangle-\langle Z,[Y,X]\rangle
+\langle X,[Z,Y]\rangle\bigr).
\tag{5.9}
\end{equation}
This is Koszul's formula.
\item[(b)] In smooth coordinates,
\begin{equation}
\Gamma^k_{ij}
=\frac12g^{kl}(\partial_ig_{jl}+\partial_jg_{il}-\partial_lg_{ij}).
\tag{5.10}
\end{equation}
\item[(c)] For a smooth local frame $(E_i)$, define $c^k_{ij}$ by
\begin{equation}
[E_i,E_j]=c^k_{ij}E_k.
\tag{5.11}
\end{equation}
Then
\begin{equation}
\Gamma^k_{ij}
=\frac12g^{kl}\left(E_ig_{jl}+E_jg_{il}-E_lg_{ij}
-g_{jm}c^m_{il}-g_{lm}c^m_{ji}+g_{im}c^m_{lj}\right).
\tag{5.12}
\end{equation}
\item[(d)] If $g$ is Riemannian and $(E_i)$ is orthonormal, then
\begin{equation}
\Gamma^k_{ij}=\frac12(c^k_{ij}-c^j_{ik}-c^i_{jk}).
\tag{5.13}
\end{equation}
\end{enumerate}

\begin{proof}
\sketch Formulas (5.9) and (5.10) follow from the theorem. Substituting
$X=E_i$, $Y=E_j$, and $Z=E_l$ into (5.9) gives
\[
\Gamma^q_{ij}g_{ql}
=\frac12\left(E_ig_{jl}+E_jg_{il}-E_lg_{ij}
-g_{jm}c^m_{il}-g_{lm}c^m_{ji}+g_{im}c^m_{lj}\right).
\]
Multiplication by $g^{kl}$ proves (5.12). If $g_{ij}=\delta_{ij}$, this
reduces to (5.13), using antisymmetry of $c^k_{ij}$ in $i,j$.
\end{proof}
\end{corollary}

Henceforth, geodesics are taken with respect to the Levi-Civita connection.
Its coordinate coefficients (5.10) are the \emph{Christoffel symbols} of $g$.

\begin{proposition}
\label{prop:levi-civita-on-euclidean-and-submanifolds}
\dcref{ch5:5.12}
\uses{prop:tangential-connection-compatible-metric,prop:tangential-connection-symmetric}

\begin{enumerate}
\item[(a)] The Levi-Civita connection of a (pseudo-)Euclidean space is its
Euclidean connection.
\item[(b)] The Levi-Civita connection of an embedded (pseudo-)Riemannian
submanifold of a (pseudo-)Euclidean space is its tangential connection
$\nabla^\top$.
\end{enumerate}

\begin{proof}
\sketch The Euclidean connection is symmetric and compatible with every
pseudo-Euclidean metric. Part (b) follows from
Propositions~\ref{prop:tangential-connection-compatible-metric}
and~\ref{prop:tangential-connection-symmetric}, together with uniqueness.
\end{proof}
\end{proposition}

\begin{proposition}[Naturality of the Levi-Civita Connection]
\label{prop:naturality-levi-civita}
\dcref{ch5:5.13}
\uses{lem:pullback-connection}

Let $(M,g)$ and $(\overline M,\overline g)$ be Riemannian or
pseudo-Riemannian manifolds, possibly with boundary, with Levi-Civita
connections $\nabla$ and $\overline\nabla$. If
$\varphi:M\to\overline M$ is an isometry, then
$\varphi^*\overline\nabla=\nabla$.

\begin{proof}
\sketch Since $\varphi$ is an isometry,
$\langle Y,Z\rangle=\langle\varphi_*Y,\varphi_*Z\rangle\circ\varphi$.
Metric compatibility of $\overline\nabla$ therefore gives
\[
X\langle Y,Z\rangle
=\langle(\varphi^*\overline\nabla)_XY,Z\rangle
+\langle Y,(\varphi^*\overline\nabla)_XZ\rangle.
\]
Naturality of brackets gives
\[
(\varphi^*\overline\nabla)_XY-(\varphi^*\overline\nabla)_YX
=(\varphi^{-1})_*[\varphi_*X,\varphi_*Y]=[X,Y].
\]
Thus the pullback connection is metric and symmetric, so uniqueness proves the
claim.
\end{proof}
\end{proposition}

\begin{corollary}[Naturality of Geodesics]
\label{cor:naturality-geodesics}
\dcref{ch5:5.14}
\uses{prop:pullback-connection-properties}

Let $\varphi:(M,g)\to(\overline M,\overline g)$ be a local isometry between
Riemannian or pseudo-Riemannian manifolds, possibly with boundary. If $\gamma$
is a geodesic in $M$, then $\varphi\circ\gamma$ is a geodesic in
$\overline M$.

\begin{proof}
\sketch Apply Proposition~4.38 locally, using naturality of the Levi-Civita
connection and the local character of the geodesic equation.
\end{proof}
\end{corollary}

\begin{proposition}
\label{prop:levi-civita-tensor-compatibility}
\dcref{ch5:5.15}

On a Riemannian or pseudo-Riemannian manifold, the connections induced by the
Levi-Civita connection on tensor bundles are compatible with the induced inner
products: for $X\in\mathfrak X(M)$ and
$F,G\in\Gamma(T^{(k,l)}TM)$,
\[
X\langle F,G\rangle
=\langle\nabla_XF,G\rangle+\langle F,\nabla_XG\rangle.
\]
\end{proposition}

\begin{proof}
\sketch Reduce by linearity to decomposable tensors
$F=\alpha_1\otimes\cdots\otimes\alpha_{k+l}$ and
$G=\beta_1\otimes\cdots\otimes\beta_{k+l}$, then apply the tensor-product
rule and the vector/covector metric-compatibility identities in (2.15).
\end{proof}

\begin{proposition}
\label{prop:riemannian-volume-form-parallel}
\dcref{ch5:5.16}
\uses{prop:tangent-bundle-admits-connection}

On an oriented Riemannian manifold $(M,g)$, the Riemannian volume form $dV_g$
is parallel for the Levi-Civita connection.

\begin{proof}
\sketch Given $p\in M$ and $v\in T_pM$, choose a curve $\gamma$ with
$\gamma(0)=p$, $\gamma'(0)=v$, and a parallel oriented orthonormal frame
$(E_1,\ldots,E_n)$ along it. Since
$dV_g(E_1,\ldots,E_n)=1$ and $D_tE_i=0$, the covariant derivative formula
gives $\nabla_vdV_g=0$.
\end{proof}
\end{proposition}

\begin{proposition}
\label{prop:musical-isomorphisms-commute}
\dcref{ch5:5.17}

Musical isomorphisms commute with total covariant differentiation. If $F$ has
a contravariant $i$th index and $\flat$ lowers that index, then
\begin{equation}
\nabla(F^\flat)=(\nabla F)^\flat.
\tag{5.14}
\end{equation}
If $G$ has a covariant $i$th index and $\sharp$ raises it, then
\begin{equation}
\nabla(G^\sharp)=(\nabla G)^\sharp.
\tag{5.15}
\end{equation}

\begin{proof}
\sketch Lowering the $i$th index is the corresponding contraction in
$F\otimes g$. Since $\nabla g=0$ and covariant differentiation commutes with
contraction,
\[
\nabla_X(F^\flat)
=\nabla_X\operatorname{tr}(F\otimes g)
=\operatorname{tr}((\nabla_XF)\otimes g)
=(\nabla_XF)^\flat.
\]
This proves (5.14). Apply it to $F=G^\sharp$ and then raise the same index to
obtain (5.15).
\end{proof}
\end{proposition}
\section{The Exponential Map}

Let $(M,g)$ be a Riemannian or pseudo-Riemannian $n$-manifold without boundary,
with its Levi-Civita connection.  For $v\in T_pM$, let $\gamma_v$ denote the
maximal geodesic with $\gamma_v(0)=p$ and $\dot\gamma_v(0)=v$.

\begin{lemma}[Rescaling Lemma]
\label{lem:rescaling-lemma}
\dcref{ch5:5.18}

For every $p \in M$, $v \in T_p M$, and $c,t \in \mathbb{R}$,
\begin{equation}
\gamma_{cv}(t) = \gamma_v(ct), \tag{5.16}
\end{equation}
whenever either side is defined.

\begin{proof}
\sketch
The assertion is immediate for $c=0$.  If $c\ne0$ and $I$ is the maximal
domain of $\gamma_v$, define $\widetilde\gamma:c^{-1}I\to M$ by
$\widetilde\gamma(t)=\gamma_v(ct)$.  Then
$\widetilde\gamma(0)=p$, $\dot{\widetilde\gamma}(0)=cv$, and
\[
\widetilde D_t\dot{\widetilde\gamma}(t)
=c^2D_t\dot\gamma_v(ct)=0.
\]
Uniqueness and maximality give $\widetilde\gamma=\gamma_{cv}$.  Applying the
same argument with scaling factor $c^{-1}$ gives the equality on the full
domain of either side.
\end{proof}
\end{lemma}

The construction below and its basic properties are summarized in
Proposition~\ref{prop:exponential-map-properties}.

Define
\[
\mathcal E=\{v\in TM:\gamma_v\text{ is defined on an interval containing }[0,1]\}
\]
and the \emph{exponential map} $\exp:\mathcal E\to M$ by
$\exp(v)=\gamma_v(1)$.  For $p\in M$, set
$\mathcal E_p=\mathcal E\cap T_pM$ and write
$\exp_p=\exp|_{\mathcal E_p}$.  The notation $\exp^G$ is reserved for the Lie
group exponential; the two exponentials agree in the bi-invariant setting of
Problem~\ref{prob:lie-group-bi-invariant-metric}, but not in general.

A subset $S$ of a vector space is \emph{star-shaped with respect to} $x\in S$
if it contains the segment joining $x$ to every point of $S$.

\begin{proposition}[Properties of the Exponential Map]
\label{prop:exponential-map-properties}
\dcref{ch5:5.19}
\uses{prob:lie-group-bi-invariant-metric,thm:geodesic-existence-uniqueness}

Let $(M, g)$ be a Riemannian or pseudo-Riemannian manifold, and let
$\exp: \mathcal{E} \to M$ be its exponential map.
\begin{enumerate}[label=(\alph*)]
\item $\mathcal{E}$ is an open subset of $TM$ containing the image of the zero
section, and each $\mathcal{E}_p \subseteq T_p M$ is star-shaped with respect
to $0$.
\item For each $v \in TM$,
\begin{equation}
\gamma_v(t) = \exp(tv) \tag{5.17}
\end{equation}
whenever either side is defined.
\item The exponential map is smooth.
\item Under the canonical identification $T_0(T_pM)\cong T_pM$,
$d(\exp_p)_0$ is the identity on $T_pM$.
\end{enumerate}

\begin{proof}
\sketch
The rescaling lemma gives $\exp(tv)=\gamma_v(t)$ and shows that every
$\mathcal E_p$ is star-shaped.  The zero section belongs to $\mathcal E$
because the corresponding geodesics are constant.

In natural coordinates $(x^i,v^i)$ on $TM$, define the geodesic vector field
by
\[
G=v^k\frac{\partial}{\partial x^k}
-\Gamma^k_{j\ell}v^jv^\ell\frac{\partial}{\partial v^k}.
\]
Its integral curves are the velocity lifts of geodesics, and its intrinsic
description is
\begin{equation}
Gf(p,v) = \frac{d}{dt}\bigg|_{t=0}
f(\gamma_v(t), \gamma_v'(t)). \tag{5.18}
\end{equation}
Thus the local formulas define a global smooth vector field.  If
$\theta:\mathcal D\to TM$ is its maximal flow, then $\mathcal D$ is open and
\[
\mathcal E=\{v\in TM:(1,v)\in\mathcal D\},
\qquad
\exp(v)=\pi\bigl(\theta(1,v)\bigr).
\]
Consequently $\mathcal E$ is open and $\exp$ is smooth.  Finally, for
$v\in T_pM$,
\[
d(\exp_p)_0(v)
=\frac{d}{dt}\bigg|_{t=0}\exp_p(tv)
=\dot\gamma_v(0)=v.
\]
\end{proof}
\end{proposition}

Corollary~\ref{cor:naturality-geodesics} implies the following naturality
property.

\begin{proposition}[Naturality of the Exponential Map]
\label{prop:naturality-exponential-map}
\dcref{ch5:5.20}
\uses{cor:naturality-geodesics}

Suppose $(M,g)$ and $(\tilde{M},\tilde{g})$ are Riemannian or
pseudo-Riemannian manifolds and $\varphi: M \to \tilde{M}$ is a local
isometry. Then for every $p \in M$, the following diagram commutes:
\[
\begin{array}{ccc}
\mathcal{E}_p & \xrightarrow{\ d\varphi_p\ } & \tilde{\mathcal{E}}_{\varphi(p)} \\
{\scriptstyle \exp_p}\downarrow & & \downarrow{\scriptstyle \exp_{\varphi(p)}} \\
M & \xrightarrow{\ \varphi\ } & \tilde{M},
\end{array}
\]
where $\mathcal E_p$ and $\widetilde{\mathcal E}_{\varphi(p)}$ are the
respective restricted exponential-map domains.
\end{proposition}

\begin{exercise}
\label{exer:prove-proposition-5-20}
\dcref{ch5:5.21}
\uses{prop:naturality-exponential-map}

Prove Proposition~\ref{prop:naturality-exponential-map}.
\end{exercise}

\begin{proposition}[Determination of Local Isometries]
\label{prop:local-isometries-completely-determined}
\dcref{ch5:5.22}
Let $(M, g)$ and $(\tilde{M}, \tilde{g})$ be Riemannian or
pseudo-Riemannian manifolds, with $M$ connected. If local isometries
$\varphi, \psi: M \to \tilde{M}$ satisfy
$\varphi(p) = \psi(p)$ and $d\varphi_p = d\psi_p$ at some $p\in M$, then
$\varphi = \psi$.
\end{proposition}

A Riemannian or pseudo-Riemannian manifold is \emph{geodesically complete} if
every maximal geodesic has domain $\mathbb R$; equivalently,
$\mathcal E=TM$.

\section{Normal Neighborhoods and Normal Coordinates}

Let $(M,g)$ be a Riemannian or pseudo-Riemannian $n$-manifold.  A
\emph{normal neighborhood} of $p\in M$ is a set $U=\exp_p(V)$ for which
$V\subseteq T_pM$ is a star-shaped neighborhood of $0$ and
$\exp_p|_V:V\to U$ is a diffeomorphism.  Such neighborhoods exist because
$d(\exp_p)_0$ is invertible.

Given an orthonormal basis $(b_i)$ of $T_pM$, let
$B:\mathbb R^n\to T_pM$ be $B(x)=x^ib_i$.  On a normal neighborhood
$U=\exp_p(V)$, the chart
\[
\varphi=B^{-1}\circ(\exp_p|_V)^{-1}
\]
defines Riemannian or pseudo-Riemannian \emph{normal coordinates} centered at
$p$.

\begin{proposition}[Uniqueness of Normal Coordinates]
\label{prop:uniqueness-normal-coordinates}
\dcref{ch5:5.23}

Let $(M,g)$ be a Riemannian or pseudo-Riemannian $n$-manifold, $p\in M$, and
$U$ a normal neighborhood of $p$.  Every normal coordinate basis is
orthonormal at $p$.  Conversely, each orthonormal basis $(b_i)$ of $T_pM$
determines a unique normal chart $(x^i)$ on $U$ satisfying
$\partial_i|_p=b_i$.  In the Riemannian case, two normal charts $(x^i)$ and
$(\tilde x^i)$ satisfy
\begin{equation}
\tilde{x}^i = A^i_j x^j \tag{5.19}
\end{equation}
for a constant matrix $(A^i_j)\in O(n)$.

\begin{proof}
\sketch
For $\varphi=B^{-1}\circ\exp_p^{-1}$,
$d\varphi_p^{-1}=d(\exp_p)_0\circ dB_0=B$, and hence
$\partial_i|_p=b_i$.  This proves existence and identifies the basis defining
the chart.  If $\widetilde\varphi=\widetilde B^{-1}\circ\exp_p^{-1}$ is
another normal chart, then
\[
\widetilde\varphi\circ\varphi^{-1}=\widetilde B^{-1}\circ B.
\]
This transition is a linear isometry, so in the Riemannian case its matrix
lies in $O(n)$; it is the identity precisely when the defining bases agree.
\end{proof}
\end{proposition}

\begin{proposition}[Properties of Normal Coordinates]
\label{prop:properties-normal-coordinates}
\dcref{ch5:5.24}
\uses{prop:uniqueness-normal-coordinates,prop:exponential-map-properties}

Let $(M,g)$ be a Riemannian or pseudo-Riemannian $n$-manifold, and let
$(U,(x^i))$ be a normal coordinate chart centered at $p\in M$.
\begin{enumerate}
\item[(a)] The coordinates of $p$ are $(0,\ldots,0)$.
\item[(b)] At $p$, $g_{ij}=\delta_{ij}$ in the Riemannian case and
$g_{ij}=\pm\delta_{ij}$ in the pseudo-Riemannian case.
\item[(c)] If $v=v^i\partial_i|_p\in T_pM$, then
\begin{equation}
\gamma_v(t)=(tv^1,\ldots,tv^n) \tag{5.20}
\end{equation}
for every $t$ in any interval containing $0$ on which $\gamma_v(t)\in U$.
\item[(d)] All Christoffel symbols vanish at $p$.
\item[(e)] All first partial derivatives of $g_{ij}$ vanish at $p$.
\end{enumerate}

\begin{proof}
\sketch
Parts (a)--(c) follow from the definition and
Propositions~\ref{prop:uniqueness-normal-coordinates}
and~\ref{prop:exponential-map-properties}.  Substituting (5.20) into the
geodesic equation and setting $t=0$ gives
\[
\Gamma^k_{ij}(p)v^iv^j=0
\]
for every $v$.  Polarization and the symmetry in $i,j$ imply
$\Gamma^k_{ij}(p)=0$.  Formula (5.2) for the Levi-Civita connection then gives
$\partial_k g_{ij}(p)=0$.
\end{proof}
\end{proposition}

Geodesics through the center that remain in a normal neighborhood are called
\emph{radial geodesics}.

\section{Tubular Neighborhoods and Fermi Coordinates}

Let $(M,g)$ be Riemannian, let $P\subseteq M$ be an embedded submanifold, and
let $\pi:NP\to P$ be its normal bundle.  With $\mathcal E\subseteq TM$ the
domain of $\exp$, set $\mathcal E_P=\mathcal E\cap NP$.  The restriction
$E=\exp|_{\mathcal E_P}:\mathcal E_P\to M$ is the \emph{normal exponential
map}.

A \emph{normal neighborhood} of $P$ is an open set $U=E(V)$ such that
$E|_V$ is a diffeomorphism and each $V\cap N_xP$ is star-shaped about $0$.  It
is a \emph{tubular neighborhood} if
\begin{equation}
V = \{(x, v) \in NP : |v|_g < \delta(x)\}, \tag{5.21}
\end{equation}
for a positive continuous function $\delta:P\to\mathbb R$.  If
$\delta\equiv\varepsilon$, it is a uniform, or $\varepsilon$-tubular,
neighborhood.

\begin{theorem}[Tubular Neighborhood Theorem]
\label{thm:tubular-neighborhood}
\dcref{ch5:5.25}

Every embedded submanifold of a Riemannian manifold $(M,g)$ has a tubular
neighborhood in $M$.  Every compact submanifold has a uniform tubular
neighborhood.

\begin{proof}
\sketch
Let $P_0=\{(x,0):x\in P\}\subseteq NP$.  At $(x,0)$, the differential of
$E$ maps $T_{(x,0)}P_0$ isomorphically onto $T_xP$ and the tangent space to
the fiber $N_xP$ isomorphically onto $N_xP$.  Since
$T_xM=T_xP\oplus N_xP$, the inverse function theorem makes $E$ a local
diffeomorphism near $P_0$.  For every $x\in P$, choose $\delta>0$ such that
this holds on
\begin{equation}
V_\delta(x) = \{(x', v') \in NP : d_g(x,x') < \delta, |v'|_g < \delta\}.
\tag{5.22}
\end{equation}
Set
\begin{equation}
\Delta(x) = \sup\{\delta \leq 1 : E \text{ is a diffeomorphism from }
V_\delta(x) \text{ to its image}\}. \tag{5.23}
\end{equation}
Then $\Delta(x)>0$, and $E$ is injective on $V_{\Delta(x)}(x)$.  The triangle
inequality gives
\begin{equation}
\Delta(x) - \Delta(x') \leq d_g(x, x') \tag{5.24}
\end{equation}
and, after interchanging $x,x'$, shows that $\Delta$ is continuous.

Define
\[
V=\{(x,v)\in NP:|v|_g<\tfrac12\Delta(x)\}.
\]
If $E(x,v)=E(x',v')$ and $\Delta(x')\leq\Delta(x)$, concatenating the two
radial geodesic segments yields
\[
d_g(x,x')\leq |v|_g+|v'|_g
<\tfrac12\Delta(x)+\tfrac12\Delta(x')\leq\Delta(x).
\]
Thus both points lie in $V_{\Delta(x)}(x)$ and are equal.  Hence $E|_V$ is an
injective local diffeomorphism and therefore a diffeomorphism onto the open
set $E(V)$, which is tubular.  If $P$ is compact, the positive continuous
function $\tfrac12\Delta$ has a positive minimum, producing a uniform
tubular neighborhood.
\end{proof}
\end{theorem}

\section{Fermi Coordinates}

Let $P$ be an embedded $p$-dimensional submanifold of a Riemannian
$n$-manifold $(M,g)$, and let $U=E(V)$ be a normal neighborhood.  Choose a
chart $(W_0,\psi)$ on $P$ and an orthonormal local frame
$(E_1,\ldots,E_{n-p})$ of $NP$ over $W_0$.  Writing
$\widetilde W_0=\psi(W_0)$, define
\[
B:\widetilde W_0\times\mathbb R^{n-p}\longrightarrow NP|_{W_0},
\qquad
B(x,v)=\bigl(q,v^jE_j|_q\bigr),
\quad q=\psi^{-1}(x).
\]
For $V_0=V\cap NP|_{W_0}$ and $U_0=E(V_0)$, the chart
$\varphi=B^{-1}\circ(E|_{V_0})^{-1}$ is equivalently
\begin{equation}
\varphi: E\left(q, v^1 E_1|_q + \cdots + v^{n-p} E_{n-p}|_q\right)
\mapsto \left(x^1(q), \ldots, x^p(q), v^1, \ldots, v^{n-p}\right). \tag{5.25}
\end{equation}
These are \emph{Fermi coordinates}.  This construction is treated in
\cite{Gra82}; see also \cite{Gra04} for tubular-neighborhood geometry.

The next result is the Fermi-coordinate counterpart of
Proposition~\ref{prop:properties-normal-coordinates}.

\begin{proposition}[Properties of Fermi Coordinates]
\label{prop:fermi-coordinates-properties}
\dcref{ch5:5.26}
\uses{prop:properties-normal-coordinates}

Let $P$ be an embedded $p$-dimensional submanifold of a Riemannian
$n$-manifold $(M,g)$, let $U$ be a normal neighborhood of $P$ in $M$, and let
$(x^1, \ldots, x^p, v^1, \ldots, v^{n-p})$ be Fermi coordinates on an open
subset $U_0\subseteq U$.  Write $x^{p+j}=v^j$ for $j=1,\ldots,n-p$.
\begin{enumerate}
\item[(a)] $P\cap U_0$ is defined by $x^{p+1}=\cdots=x^n=0$.
\item[(b)] At each $q\in P\cap U_0$,
\[
g_{ij} = g_{ji} = \begin{cases}
0, & 1 \le i \le p \text{ and } p+1 \le j \le n, \\
\delta_{ij}, & p+1 \le i, j \le n.
\end{cases}
\]
\item[(c)] For $q\in P\cap U_0$ and
$v=v^1E_1|_q+\cdots+v^{n-p}E_{n-p}|_q\in N_qP$,
\[
\gamma_v(t)=\bigl(x^1(q),\ldots,x^p(q),tv^1,\ldots,tv^{n-p}\bigr).
\]
\item[(d)] At each $q\in P\cap U_0$, $\Gamma^k_{ij}=0$ whenever
$p+1\leq i,j\leq n$.
\item[(e)] At each $q\in P\cap U_0$, $\partial_i g_{jk}(q)=0$ whenever
$p+1\leq i,j,k\leq n$.
\end{enumerate}
\end{proposition}
\section{Geodesics of the Model Spaces}

\section{Euclidean Space}

By Proposition~\ref{prop:levi-civita-on-euclidean-and-submanifolds}, the
Levi-Civita connection of Euclidean space is the Euclidean connection. Hence
constant-coefficient vector fields are parallel, the geodesics are the
constant-speed affine lines, and every Euclidean space is geodesically complete.

\section{Spheres}

A \emph{great circle} in $S^n(R)$ is an intersection $S^n(R)\cap\Pi$, where
$\Pi\subseteq\mathbb R^{n+1}$ is a two-dimensional linear subspace.

\begin{proposition}
\label{prop:maximal-geodesic-sphere}
\dcref{ch5:5.27}
\uses{cor:geodesic-submanifold-acceleration,cor:metric-connection-norm-geodesic}
A nonconstant curve on $S^n(R)$ is a maximal geodesic if and only if it is a
periodic constant-speed parametrization of a great circle. In particular,
$S^n(R)$ is geodesically complete.

\begin{proof}
\sketch For $p\in S^n(R)$,
\[
T_pS^n(R)=\{v\in\mathbb R^{n+1}:\langle p,v\rangle=0\}.
\]
Given $0\ne v\in T_pS^n(R)$, set $a=|v|/R$ and
$\widetilde v=v/a$, so $|\widetilde v|=R$. The curve
\[
\gamma(t)=(\cos at)p+(\sin at)\widetilde v
\]
lies in $S^n(R)$, satisfies $\gamma(0)=p$, $\gamma'(0)=v$, and has
\[
\gamma''(t)=-a^2\gamma(t).
\]
Its acceleration is therefore normal to the sphere, so
Corollary~\ref{cor:geodesic-submanifold-acceleration} identifies it with the
geodesic having initial data $(p,v)$. It has period $2\pi/a$, constant speed,
and image $S^n(R)\cap\operatorname{span}\{p,\widetilde v\}$. Conversely, an
orthonormal basis of any two-plane $\Pi$ gives such a parametrization of
$S^n(R)\cap\Pi$.
\end{proof}
\end{proposition}

\section{Hyperbolic Spaces}

\begin{proposition}
\label{prop:geodesics-hyperbolic-spaces}
\dcref{ch5:5.28}
A nonconstant curve in a hyperbolic space is a maximal geodesic if and only if
it is a constant-speed embedding of $\mathbb R$ with one of the following
images:
\begin{enumerate}
\item[(a)] In the hyperboloid model, $\mathbb H^n(R)\cap\Pi$ for a
two-dimensional linear subspace $\Pi\subseteq\mathbb R^{n,1}$; this is a
\emph{great hyperbola}.
\item[(b)] In the Beltrami--Klein model, the interior of a line segment whose
endpoints lie on $\partial\mathbb K^n(R)$.
\item[(c)] In the ball model, the interior of a diameter of $\mathbb B^n(R)$,
or the intersection of $\mathbb B^n(R)$ with a Euclidean circle orthogonal to
$\partial\mathbb B^n(R)$.
\item[(d)] In the half-space model, the intersection of $\mathbb U^n(R)$ with
a line parallel to the $y$-axis or with a Euclidean circle centered on
$\partial\mathbb U^n(R)$.
\end{enumerate}
\end{proposition}

\begin{theorem}
\label{thm:hyperbolic-geodesically-complete}
\uses{prop:levi-civita-on-euclidean-and-submanifolds,cor:geodesic-submanifold-acceleration}
Every hyperbolic space is geodesically complete.

\begin{proof}
\sketch Write $\overline q=q^{(n,1)}$. For $p\in\mathbb H^n(R)$,
\[
T_p\mathbb H^n(R)=\{v\in\mathbb R^{n,1}:\overline q(p,v)=0\}.
\]
For $0\ne v\in T_p\mathbb H^n(R)$, put
\[
a=\frac{|v|_{\overline q}}R=\frac{\overline q(v,v)^{1/2}}R,
\qquad \widehat v=\frac va,
\qquad
\gamma(t)=(\cosh at)p+(\sinh at)\widehat v.
\]
Then $\gamma(\mathbb R)\subseteq\mathbb H^n(R)$,
$\gamma(0)=p$, $\gamma'(0)=v$, and $\gamma''=a^2\gamma$. Thus the acceleration
is $\overline q$-normal to the hyperboloid, so $\gamma$ is the geodesic with
initial data $(p,v)$. It is a constant-speed embedding of $\mathbb R$ onto
$\mathbb H^n(R)\cap\operatorname{span}\{p,\widehat v\}$. Conversely, every
two-plane meeting $\mathbb H^n(R)$ yields such a geodesic by choosing a point
$p$ in the intersection and a nonzero vector in the plane orthogonal to $p$.
Hence the hyperboloid model is complete; naturality under the model isometries
gives completeness of the other models.

For the image classification, the isometry
\[
c\colon\mathbb H^n(R)\longrightarrow\mathbb K^n(R),
\qquad c(\xi,\tau)=\frac{R\xi}{\tau},
\]
takes the equation $\alpha_i\xi^i+\beta\tau=0$ to
$\alpha_iw^i=-\beta R$. Thus each great hyperbola maps to the intersection of
$\mathbb K^n(R)$ with a line.

In dimension two, inverse hyperbolic stereographic projection is
\[
\pi^{-1}\colon\mathbb B^2(R)\longrightarrow\mathbb H^2(R),
\qquad
\pi^{-1}(u)=\left(
\frac{2R^2u}{R^2-|u|^2},
\frac{R(R^2+|u|^2)}{R^2-|u|^2}
\right).
\]
If a great hyperbola is given by $\alpha_i\xi^i+\beta\tau=0$, the case
$\beta=0$ maps to a diameter. For $\beta\ne0$, normalize $\beta=-1$. Pullback
of $\tau=\alpha\mathbin\cdot\xi$ gives
\[
R\frac{R^2+|u|^2}{R^2-|u|^2}
=\frac{2R^2\alpha\mathbin\cdot u}{R^2-|u|^2},
\qquad
|u|^2-2R\alpha\mathbin\cdot u+R^2=0,
\]
and hence
\begin{equation}
|u-R\alpha|^2=R^2(|\alpha|^2-1). \tag{5.26}
\end{equation}
The locus meets $\mathbb B^2(R)$ only when $|\alpha|^2>1$; it is then a
circle centered at $R\alpha$ with radius $R\sqrt{|\alpha|^2-1}$. At an
intersection point $u_0\in\partial\mathbb B^2(R)$, equation (5.26) gives the
Pythagorean relation among the lengths $|u_0|=R$, $|R\alpha|$, and
$|u_0-R\alpha|$, so the circle is orthogonal to the boundary. In higher
dimensions, an orthogonal transformation in the $\xi$-variables reduces the
two-plane of a geodesic to this two-dimensional calculation; planes through
$(0,\ldots,0,R)$ give diameters.

For the upper half-space model in dimension two, use the Cayley transform
\[
\kappa(z)=iR\frac{z-iR}{z+iR}.
\]
Writing $z=x+iy$ and $\alpha=a+ib$, substitution in (5.26) reduces to
\[
(1-b)|z|^2-2aRx+(b+1)R^2=0.
\]
This is a circle centered on the $x$-axis unless $b=1$; in that case
$|\alpha|^2>1$ implies $a\ne0$, and the locus is a vertical line
$x=\text{constant}$. In higher dimensions, translate and rotate the
$x$-variables so that the initial point lies on the $y$-axis and the initial
velocity belongs to
$\operatorname{span}\{\partial/\partial x^1,\partial/\partial y\}$. The
two-dimensional calculation then gives a vertical half-line or a semicircle
centered on $y=0$, and the inverse translation and rotation preserve these
classes.
\end{proof}
\end{theorem}

\section{Euclidean and Non-Euclidean Geometries}

\subsection*{Euclidean Plane Geometry}

One rigorous formulation of Euclidean plane geometry is Hilbert's system
\cite{Hil71}; see also \cite{Gan73,Gre93} and \cite{LeeAG}. Its
primitive terms are \emph{point}, \emph{line}, \emph{lies on}, \emph{between},
and \emph{congruent}. The associated notation is fixed as follows.
\begin{itemize}
\item A line $l$ \emph{contains} $P$ when $P$ lies on $l$; a set is
\emph{collinear} when one line contains all its points.
\item For distinct $A,B$, the segment $\overline{AB}$ consists of $A$, $B$,
and the points between them, and $\overline{AB}\cong\overline{A'B'}$ denotes
segment congruence.
\item The ray $\overrightarrow{AB}$ consists of $A$, $B$, the points between
$A$ and $B$, and the points $C$ for which $B$ is between $A$ and $C$. Its
interior points are all its points except $A$.
\item For noncollinear $A,O,B$, the angle $\angle AOB$ is
$\overrightarrow{OA}\cup\overrightarrow{OB}$, and
$\angle ABC\cong\angle A'B'C'$ denotes angle congruence.
\item Points $A,B\notin l$ are on the same side of $l$ when
$\overline{AB}\cap l=\varnothing$.
\item Two lines are parallel when they have no common point.
\end{itemize}

The postulates used here are the following.
\begin{itemize}
\item \emph{Incidence:}
\begin{enumerate}
\item[(a)] Every two distinct points lie on a unique line.
\item[(b)] Every line contains at least two points, and at least three
noncollinear points exist.
\end{enumerate}
\item \emph{Order:}
\begin{enumerate}
\item[(a)] If $B$ is between $A$ and $C$, then $A,B,C$ are distinct collinear
points, and $B$ is also between $C$ and $A$.
\item[(b)] For distinct $A,C$, some $B$ has $C$ between $A$ and $B$.
\item[(c)] Of three distinct collinear points, at most one is between the
other two.
\item[(d)] If $A,B,C$ are noncollinear and a line $l$ containing none of them
meets $\overline{AB}$, then it meets $\overline{AC}$ or $\overline{BC}$.
\end{enumerate}
\item \emph{Congruence:}
\begin{enumerate}
\item[(a)] Given $A,B\in l$, a point $A'\in l'$, and a chosen ray of $l'$
from $A'$, there is a point $B'$ on that ray with
$\overline{AB}\cong\overline{A'B'}$.
\item[(b)] Any two segments congruent to the same segment are congruent to
each other.
\item[(c)] If $\overline{AB}$ and $\overline{BC}$ meet only at $B$, and
$\overline{A'B'}$ and $\overline{B'C'}$ meet only at $B'$, then
\[
\overline{AB}\cong\overline{A'B'},\quad
\overline{BC}\cong\overline{B'C'}
\quad\Longrightarrow\quad
\overline{AC}\cong\overline{A'C'}.
\]
\item[(d)] Given an angle $\angle rs$, a line $l'$, a chosen side of $l'$,
and a ray $\vec r'$ on $l'$ from $O'$, there is a unique ray $\vec s'$ whose
interior lies on that side and for which $\angle r's'\cong\angle rs$.
\item[(e)] If triangles $ABC$ and $A'B'C'$ satisfy
\[
\overline{AB}\cong\overline{A'B'},\qquad
\overline{AC}\cong\overline{A'C'},\qquad
\angle BAC\cong\angle B'A'C',
\]
then $\angle ABC\cong\angle A'B'C'$ and
$\angle ACB\cong\angle A'C'B'$.
\end{enumerate}
\item \emph{Euclidean parallel postulate:} for every line $l$ and point
$A\notin l$, exactly one line through $A$ is parallel to $l$.
\end{itemize}

An \emph{interpretation} assigns meanings to the primitive terms; it is a
\emph{model} when all the postulates hold under those meanings. The Cartesian
model assigns:
\begin{itemize}
\item points are elements of $\mathbb R^2$;
\item lines are images of maximal Euclidean geodesics;
\item $A$ lies on $l$ exactly when $A\in l$;
\item for distinct $A,B,C$, the point $B$ is between $A$ and $C$ exactly when
it lies on the geodesic segment joining them;
\item subsets $S,S'$ are congruent exactly when some Euclidean isometry
$\varphi\colon\mathbb R^2\to\mathbb R^2$ satisfies $\varphi(S)=S'$.
\end{itemize}

\begin{exercise}
\label{exer:hilbert-axioms-euclidean-plane}
\dcref{ch5:5.29}
Verify Hilbert's axioms for the Cartesian interpretation above.
\end{exercise}

Replacing the Euclidean parallel postulate by
\begin{quote}
For every line $l$ and point $A\notin l$, at least two distinct lines through
$A$ are parallel to $l$
\end{quote}
gives the hyperbolic parallel postulate. The preceding primitive terms are
interpreted on $\mathbb H^2$ using maximal geodesics and isometries of any
hyperbolic metric $g_B$; these axioms do not distinguish radii. The relevant
axioms for this model are addressed in
Problem~\ref{prob:hyperbolic-geometry-hilbert-axioms}.

The elliptic parallel postulate asserts that no two lines are parallel. Since
Hilbert's remaining axioms imply the existence of a parallel through any point
outside a line (see [Gre93]), an elliptic system must also modify other axioms.
On the round sphere $S^2$, with lines defined as images of maximal geodesics,
antipodal points lie on infinitely many lines and any two distinct lines meet
in two points. This model is called \emph{double elliptic geometry}.

\begin{example}
\label{ex:real-projective-plane-elliptic-geometry}
\dcref{ch5:5.30}
\uses{ex:real-projective-space-metric,prob:rp-n-frame-homogeneous,prob:elliptic-geometry-hilbert}
The real projective plane $\mathbb{RP}^2$ carries a frame-homogeneous metric
$g$ locally isometric to a round metric on $S^n$ (Example 2.34 and
Problem 3-2). With points given by elements of $\mathbb{RP}^2$ and lines by
images of maximal geodesics, the incidence postulates and the elliptic parallel
postulate hold by Problem~\ref{prob:elliptic-geometry-hilbert}. This is
\emph{single elliptic geometry}.
\end{example}
\section{Problems}

\begin{exercise}
\label{prob:metric-connections-difference-tensor}
\dcref{ch5:5-1}
\uses{prop:difference-tensor}

Let $(M,g)$ be Riemannian or pseudo-Riemannian, let $\nabla$ be its
Levi-Civita connection, and let $\widetilde\nabla$ be another connection on
$TM$. If $D$ is their difference tensor and
\[
D^\flat(X,Y,Z)=\langle D(X,Y),Z\rangle,
\]
show that $\widetilde\nabla$ is compatible with $g$ if and only if
\[
D^\flat(X,Y,Z)=-D^\flat(X,Z,Y)
\]
for all $X,Y,Z\in\mathfrak X(M)$. Deduce that when $\dim M\geq2$, the metric
connections form an infinite-dimensional affine space.
\end{exercise}

\begin{exercise}
\label{prob:connection-forms-compatibility}
\dcref{ch5:5-2}
\uses{prob:connection-1forms-cartan}

Let $\nabla$ be a connection on the tangent bundle of a Riemannian manifold
$(M,g)$. Show that $\nabla$ is compatible with $g$ if and only if, in every
local frame $(E_i)$, its connection $1$-forms satisfy
\[
g_{jk}\omega_i^k+g_{ik}\omega_j^k=dg_{ij}.
\]
Deduce that $(\omega_i^j)$ is skew-symmetric in every local orthonormal frame.
\end{exercise}

\begin{exercise}
\label{prob:nonsymmetric-euclidean-connection}
\dcref{ch5:5-3}
\uses{prob:connections-difference-tensor}

In standard coordinates on $\mathbb R^3$, define a connection by
\[
\Gamma_{12}^3=\Gamma_{23}^1=\Gamma_{31}^2=1,
\qquad
\Gamma_{21}^3=\Gamma_{32}^1=\Gamma_{13}^2=-1,
\]
with every other coefficient zero. Show that it is Euclidean-metric compatible
and has the same geodesics as the Euclidean connection, but is not symmetric.
\end{exercise}

\begin{exercise}
\label{prob:christoffel-symbols-surface-revolution}
\dcref{ch5:5-4}
\uses{ex:surfaces-of-revolution}

Let $C$ be an embedded smooth curve in
$H=\{(r,z):r>0\}$, and let $S_C\subseteq\mathbb R^3$ be its surface of
revolution as in Example~2.20. Let $\gamma(t)=(a(t),b(t))$ be a unit-speed
local parametrization of $C$, and let $X$ be the parametrization in (2.11).
\begin{enumerate}
\item[(a)] Compute the Christoffel symbols of the induced metric in
$(t,\theta)$ coordinates.
\item[(b)] Show that every meridian is the image of a geodesic on $S_C$.
\item[(c)] Determine necessary and sufficient conditions for a latitude circle
to be the image of a geodesic.
\end{enumerate}
\end{exercise}

\begin{exercise}
\label{prob:parallel-vector-fields}
\dcref{ch5:5-5}
\uses{ex:surfaces-of-revolution,prob:christoffel-symbols-surface-revolution}

A vector field $Y$ on an open subset of a Riemannian manifold is
\emph{parallel} if $\nabla Y=0$.
\begin{enumerate}
\item[(a)] Given $p\in\mathbb R^n$ and $v\in T_p\mathbb R^n$, prove that
there is a unique parallel vector field $Y$ on $\mathbb R^n$ satisfying
$Y_p=v$.
\item[(b)] On the unit sphere, consider
\[
X(\varphi,\theta)
=(\sin\varphi\cos\theta,\sin\varphi\sin\theta,\cos\varphi)
\]
and the associated coordinate fields $X_\theta=X_*(\partial_\theta)$ and
$X_\varphi=X_*(\partial_\varphi)$. Compute
$\nabla_{X_\theta}X_\varphi$ and $\nabla_{X_\varphi}X_\varphi$, and deduce that
$X_\varphi$ is parallel along the equator and every meridian
$\theta=\theta_0$.
\item[(c)] For $p=(1,0,0)\in S^2$, prove that no neighborhood of $p$ admits
a parallel vector field $W$ with $W_p=X_\varphi|_p$.
\item[(d)] Use (a) and (c) to prove that no neighborhood of $p$ in
$(S^2,\overline g)$ is isometric to an open subset of
$(\mathbb R^2,\overline g)$.
\end{enumerate}
\end{exercise}

\begin{exercise}
\label{prob:riemannian-submersion-connection}
\dcref{ch5:5-6}
\uses{prop:horizontal-vector-field-properties}

Let $\pi:(\widetilde M,\widetilde g)\to(M,g)$ be a Riemannian submersion, and
write $\widetilde Z$ for the horizontal lift of $Z\in\mathfrak X(M)$.
\begin{enumerate}
\item[(a)] For $X,Y\in\mathfrak X(M)$, prove
\begin{align*}
\langle\widetilde X,\widetilde Y\rangle&=\langle X,Y\rangle\circ\pi,\\
[\widetilde X,\widetilde Y]^H&=\widetilde{[X,Y]},\\
[\widetilde X,W]&\text{ is vertical whenever $W\in\mathfrak X(\widetilde M)$
is vertical.}
\end{align*}
\item[(b)] If $\widetilde\nabla$ and $\nabla$ are the Levi-Civita connections
of $\widetilde g$ and $g$, prove
\begin{equation}
\widetilde\nabla_{\widetilde X}\widetilde Y
=\widetilde{\nabla_XY}+\frac12[\widetilde X,\widetilde Y]^V.
\tag{5.27}
\end{equation}
\emph{Hint.} For a horizontal lift $\widetilde Z$ and a vertical field $W$,
compute the inner products of
$\widetilde\nabla_{\widetilde X}\widetilde Y$ with $\widetilde Z$ and $W$
using (5.9).
\end{enumerate}
\end{exercise}

\begin{exercise}
\label{prob:geodesics-product-warped-metrics}
\label{prob:geodesic-warped-product}
\dcref{ch5:5-7}
\uses{ex:warped-products}

Let $(M_1,g_1)$ and $(M_2,g_2)$ be Riemannian manifolds.
\begin{enumerate}
\item[(a)] With the product metric on $M_1\times M_2$, prove that
$\gamma(t)=(\gamma_1(t),\gamma_2(t))$ is a geodesic if and only if each
$\gamma_i$ is a geodesic in $(M_i,g_i)$.
\item[(b)] Let $f:M_1\to\mathbb R^+$ be smooth and strictly positive, and
equip $M_1\times_fM_2$ with the corresponding warped metric. For fixed
$q_0\in M_2$, prove that $\gamma(t)=(\gamma_1(t),q_0)$ is a warped-product
geodesic if and only if $\gamma_1$ is a $g_1$-geodesic.
\end{enumerate}
\end{exercise}

\begin{exercise}
\label{prob:lie-group-bi-invariant-metric}
\dcref{ch5:5-8}
\uses{prop:bi-invariant-metric-adjoint,prop:evaluation-map-lie-algebra-isomorphism,prop:properties-of-lie-group-exponential-map}

Let $G$ be a Lie group with Lie algebra $\mathfrak g$, and let $g$ be a
bi-invariant Riemannian metric whose inner product on $\mathfrak g$ is
$\langle\cdot,\cdot\rangle$.
\begin{enumerate}
\item[(a)] Prove that $\operatorname{ad}(X)$ is skew-adjoint for every
$X\in\mathfrak g$:
\[
\langle\operatorname{ad}(X)Y,Z\rangle
=-\langle Y,\operatorname{ad}(X)Z\rangle.
\]
\emph{Hint.} Differentiate
$\langle\operatorname{Ad}(\exp^G(tX))Y,
\operatorname{Ad}(\exp^G(tX))Z\rangle$ at $t=0$ and use
$\operatorname{Ad}_*=\operatorname{ad}$.
\item[(b)] For left-invariant vector fields $X,Y$ on $G$, prove
$\nabla_XY=\frac12[X,Y]$.
\item[(c)] Prove that the geodesics starting at the identity are precisely the
one-parameter subgroups. Under the canonical identification
$\mathfrak g\cong T_eG$, deduce that the restricted Riemannian exponential at
$e$ equals the Lie group exponential $\exp^G:\mathfrak g\to G$.
\item[(d)] Regard $\mathbb R^+$ as a multiplicative Lie group. Prove that
$g=r^{-2}dt^2$ is bi-invariant and that the restricted Riemannian exponential
at $1$ sends $c\,\partial/\partial t$ to $e^c$.
\end{enumerate}
\end{exercise}

\begin{exercise}
\label{prob:normal-coordinates-chart}
\dcref{ch5:5-9}

Let $(U,\varphi)$ be a coordinate chart about $p$ in a Riemannian manifold
$(M,g)$, with $\varphi(p)=0$ and $\varphi(U)$ star-shaped about $0$. Prove
that it is a normal coordinate chart if and only if
$g_{ij}(p)=\delta_{ij}$ and
\[
x^ix^j\Gamma^k_{ij}(x)=0
\]
identically on $U$.
\end{exercise}

\begin{exercise}
\label{prob:local-isometry-determined-point}
\dcref{ch5:5-10}
\uses{prop:local-isometries-completely-determined}

Prove Proposition~\ref{prop:local-isometries-completely-determined}, which
states that a local isometry is determined by its value and differential at
one point.
\end{exercise}

\begin{exercise}
\label{prob:model-manifolds-isometric-actions}
\dcref{ch5:5-11}

Consider the isometric actions of $E(n)$, $O(n+1)$, and $O^+(n,1)$ on
$(\mathbb R^n,\overline g)$, $(S^n(R),\overline g_R)$, and
$(\mathbb H^n(R),\overline g_R)$, respectively.
\begin{enumerate}
\item[(a)] Prove
\[
\operatorname{Iso}(\mathbb R^n,\overline g)=E(n),\qquad
\operatorname{Iso}(S^n(R),\overline g_R)=O(n+1),\qquad
\operatorname{Iso}(\mathbb H^n(R),\overline g_R)=O^+(n,1).
\]
\item[(b)] In each model, prove that the isotropy group at every point is
isomorphic to $O(n)$.
\item[(c)] Let $(M,g)$ be one of these three models, let $U\subseteq M$ be
connected and open, and let $\varphi:U\to M$ be a local isometry. Prove that
$\varphi$ is the restriction to $U$ of an element of
$\operatorname{Iso}(M,g)$.
\end{enumerate}
\end{exercise}

\begin{exercise}
\label{prob:lie-group-dimension-bound}
\dcref{ch5:5-12}

Let a Lie group $G$ act effectively and isometrically on a connected
$n$-dimensional Riemannian manifold $M$. Prove
$\dim G\leq n(n+1)/2$.
\end{exercise}

\begin{exercise}
\label{prob:exterior-derivative-covariant-formula}
\dcref{ch5:5-13}

Let $(M,g)$ be a Riemannian manifold.
\begin{enumerate}
\item[(a)] For every $\eta\in\Omega^1(M)$, prove
\[
d\eta(X,Y)=(\nabla_X\eta)(Y)-(\nabla_Y\eta)(X).
\]
\item[(b)] For every $\eta\in\Omega^k(M)$, prove
\[
d\eta=(-1)^k(k+1)\operatorname{Alt}(\nabla\eta),
\]
where $\operatorname{Alt}$ is the alternation operator in (B.9).
\emph{Hint.} Compute at each point in normal coordinates and use the
coordinate independence of both sides.
\end{enumerate}
\end{exercise}

\begin{exercise}
\label{prob:divergence-laplacian-formulas}
\dcref{ch5:5-14}

Let $(M,g)$ be a Riemannian manifold, and let $\operatorname{div}$ and
$\Delta$ denote its divergence and Laplace operators.
\begin{enumerate}
\item[(a)] For $X\in\mathfrak X(M)$, prove
\[
\operatorname{div}X=\operatorname{tr}(\nabla X),
\qquad
\operatorname{div}X=X^i_{;i}
\]
when $X=X^iE_i$ in a local frame. \emph{Hint.} It suffices to calculate at
the origin of normal coordinates.
\item[(b)] For $u\in C^\infty(M)$, prove
\[
\Delta u=\operatorname{tr}_g(\nabla^2u)
=g^{ij}u_{;ij}=u_{;i}^{\ i}
\]
in every local frame.
\end{enumerate}
\end{exercise}

\begin{exercise}
\label{prob:eigenfunction-laplacian-inequality}
\dcref{ch5:5-15}
\uses{prob:eigenvalues,prob:harmonic-functions}

Let $(M,g)$ be a compact Riemannian $n$-manifold without boundary, and suppose
$u\in C^\infty(M)$ satisfies $-\Delta u=\lambda u$ for a constant $\lambda$.
Prove
\[
\lambda\int_M|\operatorname{grad}u|^2\,dV_g
\leq n\int_M|\nabla^2u|^2\,dV_g.
\]
\emph{Hint.} Consider $\nabla^2u-\frac1n(\Delta u)g$ and use one of Green's
identities from Problem~2-23.
\end{exercise}

\begin{exercise}
\label{prob:divergence-tensor-fields}
\dcref{ch5:5-16}
\uses{prob:divergence-laplacian-formulas,prob:divergence-theorem}

For a smooth tensor field $F$ of any covariant, contravariant, or mixed type on
a Riemannian manifold, define
\[
\operatorname{div}F=\operatorname{tr}_g(\nabla F),
\]
where the trace contracts the last two indices. If $F$ is purely
contravariant, ordinary trace may be used because the penultimate index of
$\nabla F$ is contravariant.

Let $F$ be a smooth covariant $k$-tensor and $G$ a smooth covariant
$(k+1)$-tensor on a compact Riemannian manifold $(M,g)$ with boundary. Prove
\[
\int_M\langle\nabla F,G\rangle\,dV_g
=\int_{\partial M}\langle F\otimes N^\flat,G\rangle\,d\widehat V_g
-\int_M\langle F,\operatorname{div}G\rangle\,dV_g,
\]
where $\widehat g$ is the metric induced on $\partial M$. Equivalently,
\[
\int_M F_{i_1\ldots i_k;j}G^{i_1\ldots i_kj}\,dV_g
=\int_{\partial M}F_{i_1\ldots i_k}G^{i_1\ldots i_kj}N_j\,dV_{\widehat g}
-\int_MF_{i_1\ldots i_k}G^{i_1\ldots i_kj}_{\ \ \ \ \ \ \ ;j}\,dV_g.
\]
\end{exercise}

\begin{exercise}
\label{prob:tubular-neighborhood-normal-bundle}
\dcref{ch5:5-17}
\uses{thm:tubular-neighborhood}

Let $P$ be an embedded submanifold of a Riemannian manifold $(M,g)$. Prove
that $P$ has a tubular neighborhood diffeomorphic to the total space of its
normal bundle $NP$, with the zero section mapped to $P$.
\emph{Hint.} First show that the function $\delta$ in (5.21) can be chosen
smooth.
\end{exercise}

\begin{exercise}
\label{prob:fermi-coordinates}
\dcref{ch5:5-18}
\uses{prop:fermi-coordinates-properties}

Prove Proposition~\ref{prop:fermi-coordinates-properties} on the properties of
Fermi coordinates.
\end{exercise}

\begin{exercise}
\label{prob:hyperbolic-geometry-hilbert-axioms}
\dcref{ch5:5-19}

Using $(\mathbb H^2,\widetilde g)$, construct an interpretation of Hilbert's
axioms with the hyperbolic parallel postulate in place of the Euclidean one.
Prove the incidence postulates, congruence postulate (e), and hyperbolic
parallel postulate in this geometry.
\end{exercise}

\begin{exercise}
\label{prob:elliptic-geometry-hilbert}
\dcref{ch5:5-20}
\uses{ex:real-projective-plane-elliptic-geometry}

In single elliptic geometry from
Example~\ref{ex:real-projective-plane-elliptic-geometry}, take points to be
elements of $\mathbb RP^2$ and lines to be images of maximal geodesics. Prove
Hilbert's incidence postulates and the elliptic parallel postulate.
\end{exercise}

\begin{exercise}
\label{prob:orthonormal-frame-zero-covariant-derivative}
\dcref{ch5:5-21}

Let $(M,g)$ be Riemannian or pseudo-Riemannian, let $p\in M$, and let
$(b_1,\ldots,b_n)$ be an orthonormal basis of $T_pM$. Prove that a neighborhood
of $p$ admits a smooth orthonormal frame $(E_i)$ satisfying
\[
E_i|_p=b_i,
\qquad
(\nabla E_i)_p=0
\]
for every $i$.
\end{exercise}

\begin{exercise}[Killing vector field]
\label{prob:killing-vector-field}
\dcref{ch5:5-22}
\uses{prop:invariant-tensor-characterization,prop:lie-derivative-tensor-product-rule}

A smooth vector field $X$ on a Riemannian manifold is a \emph{Killing vector
field} if $\mathcal L_Xg=0$, equivalently if $g$ is invariant under the flow
of $X$. Prove that $X$ is Killing if and only if the covariant $2$-tensor
$(\nabla X)^\flat$ is antisymmetric. \emph{Hint.} Use Proposition~B.9.
\end{exercise}

\begin{exercise}
\label{prob:holonomy-group}
\dcref{ch5:5-23}

Let $(M,g)$ be a connected Riemannian manifold and $p\in M$. For each
admissible loop $\gamma:[a,b]\to M$ based at $p$, let
$P^\gamma=P^\gamma_{ab}:T_pM\to T_pM$ be parallel transport, and define
\[
\operatorname{Hol}(p)
=\{P^\gamma:\gamma\text{ is an admissible loop based at }p\}
\subseteq\operatorname{GL}(T_pM).
\]
\begin{enumerate}
\item[(a)] Prove that $\operatorname{Hol}(p)$ is a subgroup of $O(T_pM)$,
called the \emph{holonomy group} at $p$.
\item[(b)] Let $\operatorname{Hol}^0(p)$ consist of transports along loops
path-homotopic to the constant loop. Prove that it is a normal subgroup of
$\operatorname{Hol}(p)$, called the \emph{restricted holonomy group}.
\item[(c)] For $p,q\in M$, construct an isomorphism
$\operatorname{GL}(T_pM)\to\operatorname{GL}(T_qM)$ carrying
$\operatorname{Hol}(p)$ to $\operatorname{Hol}(q)$.
\item[(d)] Prove that $M$ is orientable if and only if
$\operatorname{Hol}(p)\subseteq\operatorname{SO}(T_pM)$ for some $p\in M$.
\item[(e)] Prove that $g$ is flat if and only if
$\operatorname{Hol}^0(p)$ is trivial for some $p\in M$.
\end{enumerate}
\end{exercise}
